\answer{ex:12:1}{$0.60$ (sin correlación sería $0.18$), límite $\sqrt{c/(rT)}\approx0.58$: cien neuronas solo distinguen “sí/no”.}
\answer{ex:12:2}{(a)~$10^8$ espigas, $\approx10\,\text{mJ}$: el $0.3\,\%$ de los $3$~J de todo el cerebro. (b)~$3$~J, $\approx300$ veces más.}
\answer{ex:12:3}{(a)~$\theta_A\in[2,3)$, $\theta_B\in[1,2)$; la otra figura da como máximo $2$ y $1$. (b)~$10\cdot96^2\approx9.2\cdot10^4$.}
\answer{ex:12:4}{(a)~$T_d=0.85\,\tau_a=0.34\,\text{s}$. (b)~$T_d\propto\tau_a$ y no depende de $I$: $\tau_a\approx3.5\,\text{s}$ (simulación: $3.1\,\text{s}$).}
\answer{ex:12:5}{$p\approx0.17$. La fracción cae despacio: $0.90$, $0.88$, $0.86$, $0.84$, $0.82$, $0.79$ con $N=8$, $12$, $24$, $48$, $96$, $192$.}
\answer{ex:12:6}{Las diez corridas no tienen errores con $\tau_{\text{tr}}=20$, $50$ y $200\ms$ (con $30\ms$ falló una corrida de diez); con $5$--$10\ms$, alrededor de $0.5$; con $0.5$--$2\,\text{s}$ la calidad cae a $0.86$. Por arriba: la traza debe extinguirse durante la pausa, $e^{-0.3\,\text{s}/\tau_{\text{tr}}}\ll1$.}
\answer{ex:12:7}{Un solo retardo no basta: la ventana de $\pm5\ms$ no cubre $\Delta x/v$ de $5$ a $20\ms$. Dos ramas de inhibición con $5<d_1\le10\ms$ y $15\le d_2\le d_1+10\ms$.}
\answer{ex:12:8}{\texttt{two\_stage}: $\Delta=2$; un empate con una inversión, cambio de respuesta a partir de dos. Contador de unos: una inversión; $0.57$ con $p=0.1$.}
